\originalchapter{8.333作业3：动力学}
\psource{Mehran Kardar，8.333 Fall 2013，PS3，截止2013年10月18日。}
\begin{exercise}[原题1(a)--(b)]\label{p:333-3-1}
(a)证明沿Hamilton轨道$\dot O=\{O,H\}$。(b)若任意可测光滑观测$O$的系综平均$\avg{\{O,H\}}=0$，证明$\{H,\rho\}=0$，注明边界条件。
\end{exercise}
\begin{solution}
链式法则与Hamilton方程给(a)。对任意紧支撑$O$或有消失表面项的周期/衰减边界，分部积分$\int\rho\{O,H\}=\int O\{H,\rho\}$。左侧全为0，测试函数任意使$\{H,\rho\}=0$按分布意义；若$\rho$可微则逐点成立。只对一个$O$或有限几个矩为0不够推出全密度平稳。
\end{solution}
\begin{exercise}[原题2(a)--(f)]\label{p:333-3-2}
局域三维Maxwell分布$f_1^0=n(2\pi m\kb T)^{-3/2}\exp[-(p-mu)^2/(2m\kb T)]$。(a)求压力张量$P_{ij}=mn\avg{c_ic_j}$与热流$h_i=nm\avg{c_ic^2}/2$，$c=p/m-u$。(b)代入粒子、动量与能量守恒求$n,u,T$的零阶流体方程。(c)证明$D_t\log(nT^{-3/2})=0$。(d)积分动量写$H_0=\int f_1^0\log f_1^0$。(e)求$\dot H_0$。(f)解释平衡趋近。
\end{exercise}
\begin{solution}
Gaussian二阶矩给$P_{ij}=n\kb T\delta_{ij}$，奇矩$h_i=0$。故$D_tn=-n\nabla\cdot u$、$mnD_tu=-\nabla(n\kb T)$、$D_tT=-(2/3)T\nabla\cdot u$。前后两对数导数相消得(c)。积分$\avg{mc^2/(2\kb T)}=3/2$，
\begin{equation}
H_0=\int\dd^3r\,n\left[\log\frac n{(2\pi m\kb T)^{3/2}}-\frac32\right].
\end{equation}
常相空间单位仅加与$N$有关常数。由(c)方括号沿流线守恒，再以连续性求$\partial_t(n\mathcal A)+\nabla\cdot(nu\mathcal A)=0$，无边界流时$\dot H_0=0$。零阶Euler不产生耗散，局域Maxwell但非均匀态可持续运动；趋平衡需要梯度修正的黏性、热导或其他不可逆机制。
\end{solution}
\begin{exercise}[原题3(a)--(c)]\label{p:333-3-3}
两板距$w$，下板静止、上板沿$x$速度$u$。零阶分布为均匀$n,T$、局域流速$u_x=\alpha y$的Maxwell，$\alpha=u/w$。(a)用单碰撞时间$\tau_\times$求一阶$f_1=f_0-\tau_\times(v\cdot\nabla)f_0$。(b)求穿过$y$面$x$动量流$\Pi_{xy}$。(c)求板受力与黏度。
\end{exercise}
\begin{solution}
令$c_x=p_x/m-\alpha y,c_y=p_y/m$，$\partial_yf_0=(m\alpha c_x/\kb T)f_0$，所以$\delta f=-\tau_\times m\alpha c_xc_yf_0/(\kb T)$。$\Pi_{xy}=\int mc_xc_y\delta f\dd p=-n\kb T\tau_\times\alpha$，因$\avg{c_x^2c_y^2}=(\kb T/m)^2$。下板所受沿流方向的力密度$F_x=-\Pi_{xy}$，上板反向，$\eta=n\kb T\tau_\times$，与$y$无关。要求$|\alpha|\tau_\times\ll1$且板间有流体窗口$w\gg\ell$。
\end{solution}
\begin{exercise}[原题4(a)--(e)]\label{p:333-3-4}
固定两态原子$a_0,a_1$相差$\varepsilon$，与$\omega=\varepsilon/\hbar$光子作用。自发$a_1\to a_0+\gamma$率$\sigma_{sp}$，吸收截面$\sigma_{ad}$、受激发射$\sigma_{st}$均角无关。(a)求$n_1/n_0$。(b)写光子密度$f(r,p,t)$动力学。(c)求$f_{eq}$。(d)用$ f_{eq}=h^{-3}/(\e^{\beta\hbar\omega}-1)$求三截面关系。(e)介质占$x>0$，沿$+x$入射光的非平衡通量如何衰减？求穿透长。
\end{exercise}
\begin{solution}
(a)等简并两态$n_1/n_0=\e^{-\beta\varepsilon}$。选原题相空间密度归一，自发源每态为$n_1\sigma_{sp}$，吸收与受激的速度/几何常数吸入其率系数，$\partial_tf+c\hat p\cdot\nabla f=n_1\sigma_{sp}+(n_1\sigma_{st}-n_0\sigma_{ad})f$。稳态$ f_{eq}=n_1\sigma_{sp}/(n_0\sigma_{ad}-n_1\sigma_{st})=\sigma_{sp}/[\e^{\beta\varepsilon}\sigma_{ad}-\sigma_{st}]$。

(d)对所有温度匹配Planck式，需$\sigma_{ad}=\sigma_{st}$且$\sigma_{sp}=\sigma_{ad}/h^3$于上述密度归一；若截面用单位实际光通量，$c$与态密度常数须恢复，两Einstein B系数相等的内容不变。(e)沿入射方向$c\partial_x(f-f_{eq})=-(n_0\sigma_{ad}-n_1\sigma_{st})(f-f_{eq})$，所以穿透长$\ell_\gamma=c/(n_0\sigma_{ad}-n_1\sigma_{st})$，正温平衡$n_0>n_1$使衰减。若$n_1>n_0$人口反转，系数负，变增益而非正穿透长。
\end{solution}
\begin{exercise}[原题5(a)--(b)]\label{p:333-3-5}
外势$U(q)$中经典气体。(a)在归一与固定平均总能下极小一体$H=N\int\rho_1\log\rho_1\dd p\dd q$，求平衡密度。(b)两种质量$m_a,m_b$的气体共享固定总能，求两密度并说明每粒子平动能可定义共同温度。原题声称平稳解的必要性须检查其他守恒量。
\end{exercise}
\begin{solution}
(a)乘子变分$\log\rho_1+1+\lambda+\beta[p^2/(2m)+U]=0$，得$\rho_1=Z_1^{-1}\e^{-\beta[p^2/(2m)+U]}$。严格凸$\rho\log\rho$使给定这些约束下极小唯一。(b)各组归一乘子不同，能量乘子$\beta$相同，所以$\rho_\alpha\propto\e^{-\beta[p^2/(2m_\alpha)+U_\alpha]}$，每粒子平动能$3/(2\beta)=3\kb T/2$相同。质量改变速度方差，不改共同能量温标。

H定理的零碰撞态可含流速；在周期系统守恒总动量或有整体旋转等额外约束时，平稳态还可为广义canonical分布。题目所列零流唯一必要性只在无额外宏观守恒、固定静止边界或相应约束已排除时成立，不能将熵极小解当任意平稳解的无条件必要式。
\end{solution}
\begin{exercise}[原题6(a)--(c)，选做]\label{p:333-3-6}
经典电子气被相对真空深度$\varphi$的势阱束缚。(a)平坦表面逃逸电子内外速度关系。(b)温度$T$的速度分布。(c)密度$n$时热发射电流密度。
\end{exercise}
\begin{solution}
切向速度守恒。法向速度满足$v_z'^2=v_z^2-2\varphi/m$，因此仅$v_z>\sqrt{2\varphi/m}$的电子向外逃逸。设Maxwell速度密度为
\begin{equation}
P(v)=\left(\frac{m}{2\pi\kb T}\right)^{3/2}\e^{-mv^2/(2\kb T)}.
\end{equation}
电流大小由通量积分求出：
\begin{equation}
j=en\int_{v_z>v_0}v_zP(v)\dd^3v
=en\sqrt{\frac{\kb T}{2\pi m}}\e^{-\varphi/\kb T}.
\end{equation}
用入射到表面的速度做通量积分，不能用外部变小的速度却不换密度Jacobian。题目是经典模型；真实简并金属Richardson--Dushman前因子需量子态填充。
\end{solution}
